\chapter{Introducción}
\label{chap:introduccion}

\section{Objetivos}
\label{sec:objetivos}

\begin{itemize}
  \item Estudiar los conceptos básicos sobre el análisis y diseño de filtros activos.
  \item Aplicar las técnicas explicadas en clase para el cálculo de filtros activos.
  \item Implementar filtros activos.
\end{itemize}

\section{Descripción del trabajo práctico}
\label{sec:descripcion}

El trabajo práctico se divide en dos actividades. La primera (capítulo~\ref{chap:senal-campo}) consiste en el diseño
e implementación de un módulo de filtrado activo para una señal de campo de $\SI{4}{\milli\ampere}$ a
$\SI{20}{\milli\ampere}$, compuesto por un amplificador de instrumentación, un filtro pasa bajos Sallen-Key de
segundo orden y un convertidor tensión-corriente, todo alimentado con una fuente simple de $\SI{24}{\volt}$.

La segunda actividad (capítulos~\ref{chap:pasa-bajos} a~\ref{chap:pasa-banda}) consiste en el diseño de tres filtros
con los parámetros asignados al grupo, resumidos en la tabla~\ref{tab:asignacion}. Para cada uno se realizan los
cálculos, las simulaciones con elementos ideales y con valores comerciales, y para dos de ellos la implementación
física y el relevamiento de la respuesta en frecuencia.

\begin{table}[!ht]
  \centering
  \resizebox{\textwidth}{!}{\begin{tabular}{|l|c|c|c|c|c|c|}
    \hline
    Filtro & Orden & $A_{max}$ [dB] & $f_p$ / $f_o$ [kHz] & $BW$ [kHz] & Ganancia [V/V] & Aproximación / Topología \\
    \hline
    Pasa bajos & 4 & 1,5  & 5,5 & --  & 4    & Butterworth / Sallen-Key \\
    Pasa altos & 4 & 2,25 & 33  & --  & 12,5 & Butterworth / Sallen-Key \\
    Pasa banda & 2 & 0,1  & 28  & 1,7 & 9    & Butterworth / Sallen-Key \\
    \hline
  \end{tabular}}
  \caption{Parámetros asignados al Grupo 2.}
  \label{tab:asignacion}
\end{table}

% TODO: indicar cuáles dos filtros se implementan físicamente y cuál queda para el TP de filtros digitales.
